\chapter{Preces}
\label{preces}

Kyrieleison. Kyrieleison. Kyrieleison.\\
Christeleison. Christeleison. Christeleison.\\
Kyrieleison. Kyrieleison. Kyrieleison.\edissue{compare other contemporary local sources how kyrie was most usually written}
\\
Pater noster. Et ne nos.\edissue{are these incipits of recited portions, or is the rest said silently?}\\
In pace in idipsum. Dormiam et requiescam.\\
Credo in Deum. Carnis resurrectionem. Et vitam.
\edissue{are these incipits of recited portions, or is the rest said silently? Clarify the execution, provide full texts.}
\\
Benedicamus Patrem et Filium cum Sancto Spiritu.
Laudemus et superexaltemus sum in saecula.\\
Benedictus es, Domine, in firmamento coeli.
Et laudabilis et gloriosus et superexaltatus in saecula.\\
Benedicat nos omnipotens Dominus. Amen.\edissue{uncertain}
\\
Dignare Domine nocte ista.
Sine peccato nos custodire.\\
Miserere nostri, Domine, miserere nostri.\\
Fiat misericordia tua, Domine, super nos.
Quemadmodum speravimus in te.\\
Sacerdotes tui induantur iustitia.\edissue{verify exact spelling, this one differs from now standard Vulgate}
Et sancti tui exsultent.

\noindent\rubrica{\editorial{Versus proprius tempori aut festo hic adnectitur.}}

\noindent{}Domine exaudi orationem meam.
Et clamor meus \editorial{ad te veniat}.\\
Dominus vobiscum.\edissue{This dialogue is probably a standard introduction for the prayer, not part of the preces as such.}
\editorial{Et cum spiritu tuo.}

\footcite[1v]{xiv_a_19}
\edissue{Try to find in printed breviaries -- so far no luck}
